\Section{Step-by-Step Example}
\label{sec:example}

We illustrate the workflow through a number of examples; both they and the tool are publicly available \cite{InferenceExamples, MoveFlow}.

\SubSection{Higher-Order Vector Search}
\label{sec:find}

We start with a simple function that finds an element in a vector matching a predicate:

\begin{MoveBox}
  fun find<T>(v: &vector<T>,
              pred: |&T|bool has copy + drop): u64 {
    let i = 0;
    let n = v.length();
    while (i < n) {
      if (pred(&v[i])) return i;
      i = i + 1;
    }
    n
  }
\end{MoveBox}

We ask Claude Code (\CC) to infer specs for this function; this loads a skill that creates a list of tasks:

\begin{claude}
> infer specs for the functions in this package
● Skill(move-flow:move-inf)
  ⎿ Successfully loaded skill
✽ Zigzagging…
  ⎿ ◻ Synthesize loop invariants
    ◻ Infer weakest preconditions
    ◻ Simplify inferred specs
    ◻ Full-scope verification run
    ◻ Fix logical errors
    ◻ Resolve timeouts
\end{claude}

\noindent The tasks are listed in the suggested execution order but operate on independent snapshots of the package; \CC may interleave them with edits that the user can monitor or override.

\Paragraph{Synthesizing Loop Invariants}
For this task, \CC infers loop invariants for loops that do not yet have any.
The skill defines Move-specific guidelines for this. For !find!, the resulting code is:

\begin{MoveBox}
  public fun find<T>(..)
    ..
    while (i < n) {
      if (pred(vector::borrow(v, i))) return i;
      i = i + 1;
    } spec {
      invariant i <= n;
      invariant n == len(v);
      invariant forall j in 0..i: !result_of<pred>(v[j]);
      invariant forall j in 0..i: !aborts_of<pred>(v[j]);
    };
    n
  }
\end{MoveBox}

\noindent Notice the behavioral predicates !result_of! and !aborts_of! accessing the specification of another function --- here the runtime value !pred! \footnote{Specifications with function values are described in \cite{GrieskampEtAl26HO}}. Behavioral predicates are key for inference: we can refer to what other functions compute while abstracting away how, localizing inference to intra-function analysis.

\Paragraph{Inferring Weakest Preconditions}
For this task, \CC calls \MVP's \WP algorithm via its MCP service, producing:

\begin{MoveBox}
  spec find {
    pragma opaque = true;
    ensures
      (forall x: u64: x < len(v)
                      ==> !result_of<pred>(v[x]))
        ==> result == len(v);
    ensures
      forall y: u64: (forall x: u64:
                 x < y ==> !result_of<pred>(v[x]))
        && (y < len(v) && result_of<pred>(v[y]))
        ==> result == y;
    // aborts_if omitted
  }
\end{MoveBox}
%    aborts_if  exists y: u64: (forall x: u64: x < y ==> !result_of<pred>(v[x])) && ((forall x: u64: x < y ==> !aborts_of<pred>(v[x])) && (y < len(v) && aborts_of<pred>(v[y])));
%    aborts_if  exists y: u64: (forall x: u64: x < y ==> !result_of<pred>(v[x])) && ((forall x: u64: x < y ==> !aborts_of<pred>(v[x])) && (y < len(v) && !in_range(v, y)));

\noindent The !aborts_if! clauses have been omitted for brevity.


\Paragraph{Simplification of Inferred Specs}
\WP derives semantically correct but not always ideal specifications.
With loop invariants from \CC and \WP-derived specs from the Move Prover, we now ask \CC to simplify the result.

\begin{MoveBox}
  spec find {
    pragma opaque = true;
    ensures result <= len(v);
    ensures forall j: u64 where j < result:
              !result_of<pred>(v[j]);
    ensures result < len(v)
              ==> result_of<pred>(v[result]);
    // aborts_if omitted
  }
\end{MoveBox}

%    aborts_if
%        exists j: u64 where j < len(v):
%          aborts_of<pred>(v[j])
%          && (forall k: u64 where k < j:
%                !result_of<pred>(v[k]));

\noindent This specification is more compact than the \WP baseline. In general, \CC frequently rewrites \WP-generated specs into more compact and readable form, despite significant term simplification already in the \WP engine.

\Paragraph{Verification}
The remaining three tasks are verification-driven: a fast full-scope run surfaces logical errors, which are fixed iteratively; timeouts are then resolved by lemmas, proof hints, and helper functions (Sec.~\ref{sec:skills}). In this example, verification immediately succeeds.

\SubSection{Exponentiation}
\label{sec:pow}

Next, an example in the arithmetic domain.

\begin{MoveBox}
  fun pow(base: u64, exp: u64): u64 {
    let result = 1;
    let i = 0;
    while (i < exp) {
      result = result * base;
      i = i + 1;
    }
    result
  }
\end{MoveBox}

In the first step, \CC derives loop invariants, introducing a recursive helper function:

\begin{MoveBox}
  fun pow(base: u64, exp: u64): u64 {
    ..
    while (i < exp) {
      ..
    } spec {
      invariant i <= exp;
      invariant result == pow_spec(base, i);
    }
    result
  }
  spec fun pow_spec(base: u64, exp: u64): u64 {
    if (exp == 0) 1
    else base * pow_spec(base, exp - 1)
  }
\end{MoveBox}

\WP analysis then adds:

\begin{MoveBox}
  spec pow {
    ensures result == pow_spec(base, exp);
    aborts_if exp > 0
      && pow_spec(base, exp - 1) * base > MAX_U64;
  }
\end{MoveBox}

 \noindent \CC-based simplification does not find anything to improve here. However, as is to be expected with non-linear arithmetic, verification times out.

Following the timeout-resolution instructions, \CC introduces a monotonicity lemma via Move's proof-hint system.
The lemma is required for the !aborts_if! condition: if multiplication aborts with overflow earlier in the loop, then by monotonicity it also aborts later.
Below, !num! is Move's unbounded-precision signed integer type, available in specifications:

\begin{MoveBox}
  spec pow {
    .. // function spec above
  } proof {
    // Apply lemma with trigger
    forall x: num, y: num
      {pow_spec(base, x), pow_spec(base, y)}
      apply pow_mono(base, x, y);
  }
  spec lemma pow_mono(base: num, x: num, y: num) {
    requires base >= 1 && 0 <= x && x <= y;
    ensures pow_spec(base, x) <= pow_spec(base, y);
  } proof {
      if (x < y) {
          assert pow_spec(base, y - 1)
                   <= pow_spec(base, y);
          apply pow_mono(base, x, y - 1);
      }
  }
\end{MoveBox}

\SubSection{State-Dependent Inference}

As a final example, we look at a function that modifies global storage at multiple addresses in sequence, illustrating state-dependent specifications.

\begin{MoveBox}
  struct Balance has key { v: u64 }
  fun split_balance(
        dst1: &signer, dst2: &signer,
        src: address): u64 {
    let Balance { v } = move_from<Balance>(src);
    let half = v / 2;
    move_to(dst1, Balance { v: half });
    move_to(dst2, Balance { v: v - half });
    half
  }
\end{MoveBox}

\noindent After running spec inference and simplification, we obtain:

\begin{MoveBox}
  spec split_balance(
         dst1: &signer, dst2: &signer,
         src: address): u64 {
    pragma opaque = true;
    modifies Balance[src];
    modifies Balance[signer::address_of(dst1)];
    modifies Balance[signer::address_of(dst2)];
    ensures result == old(Balance[src]).v / 2;
    ensures ..S1 |~ remove<Balance>(src);
    ensures S1..S2 |~
      publish<Balance>(signer::address_of(dst1),
                       Balance{v: old(Balance[src]).v / 2});
    ensures S2.. |~
      publish<Balance>(signer::address_of(dst2),
                       Balance{v: old(Balance[src]).v
                                  - old(Balance[src]).v / 2});
    aborts_if !exists<Balance>(src);
    aborts_if S1 |~
      exists<Balance>(signer::address_of(dst1));
    aborts_if S2 |~
      exists<Balance>(signer::address_of(dst2));
  }
\end{MoveBox}

\noindent !S1! and !S2! are state labels naming intermediate memory states --- here, the memory between the !move_from! and the first !move_to!, and between the two !move_to!s, respectively.  The modality !..S1 |~ e!, !S1..S2 |~ e!, !S2.. |~ e! evaluates !e! over the named memory range, while !S |~ e! evaluates at the single state !S!.  The global-memory operators are !publish<R>(addr, val)!, which asserts that resource !R! at !addr! equals !val!, !remove<R>(addr)!, which asserts that no resource of type !R! is present at !addr!, and !exists<R>(addr)!, the standard Move resource existence predicate.
